% hw3-predicate-logic.tex (slides)

\begin{frame}[allowframebreaks]{作业 3：谓词逻辑}

\begin{problem}[一阶谓词逻辑: 形式化描述与推理]
  请使用一阶谓词逻辑公式描述以下两个定义, 并说明两者之间是否有强弱之分：
  (1) 点态连续; (2) 一致连续。
\end{problem}

\pause
\begin{solution}
  取论域为 $\mathbb{R}$，分别可形式化为
  \begin{align*}
    &\forall x\; \forall \epsilon>0\; \exists \delta>0\; \forall y\; (|x-y|<\delta \to |f(x)-f(y)|<\epsilon), \\
    &\forall \epsilon>0\; \exists \delta>0\; \forall x,y\; (|x-y|<\delta \to |f(x)-f(y)|<\epsilon).
  \end{align*}
  从量词排列可知一致连续严格强于点态连续。
\end{solution}

\end{frame}
